\documentclass[11pt]{article}
\usepackage[a4paper,margin=22mm]{geometry}
\usepackage{fontspec}
\setmainfont{DejaVu Serif}
\setsansfont{DejaVu Sans}
\usepackage{microtype}
\usepackage{xcolor}
\usepackage{hyperref}
\usepackage{enumitem}
\usepackage{parskip}
\definecolor{Accent}{HTML}{971D32}
\definecolor{Muted}{HTML}{666666}
\hypersetup{colorlinks=true,urlcolor=Accent,linkcolor=Accent}
\setlist{leftmargin=1.6em,itemsep=0.3em,topsep=0.35em}
\setcounter{tocdepth}{2}
\renewcommand{\contentsname}{Contents}
\newcommand{\sectionrule}{\par\vspace{-0.5em}\noindent\textcolor{Accent}{\rule{\linewidth}{0.6pt}}\par}
\begin{document}

\begin{center}
{\sffamily\bfseries\color{Accent} TOEFL WORD OF THE DAY\par}
\vspace{8mm}
{\fontsize{42}{48}\selectfont\bfseries delimit\par}
\vspace{2mm}
{\Large\sffamily\color{Accent} /diːˈlɪmɪt/\par}
\vspace{2mm}
{\small\sffamily Local audio phrase: delimit\par}
{\small\sffamily Audio file: audios/delimit.mp3\par}
\vspace{2mm}
{\large\sffamily\color{Muted} transitive verb\par}
\vspace{5mm}
{\sffamily mark physical boundaries \textbullet\ define the scope or limits\par}
\vspace{6mm}
{\large To delimit something is to establish where it begins and ends, either physically or conceptually.\par}
\vspace{6mm}
\begin{minipage}{0.84\textwidth}
\itshape “The researchers carefully delimited the scope of their investigation.”
\end{minipage}\par
\vspace{10mm}
\textbf{Date:} August 2nd 2026\quad\textbullet\quad \textbf{Level:} B1--mid-B2 TOEFL
\vfill
Created and compiled by \textbf{Amir Kooshky}\par
\vspace{2mm}
\href{https://t.me/KooshkyTOEFL}{Telegram Channel}\quad\textbullet\quad
\href{https://www.instagram.com/kooshkytoefl}{Instagram}\quad\textbullet\quad
\href{https://t.me/KooshkyTOEFL\_pv}{Personal Telegram}
\end{center}
\newpage
\tableofcontents
\newpage

\section{Meanings and usage}
\sectionrule
\subsection{Meaning 1: Mark physical boundaries}
\textbf{Part of speech:} transitive verb

\textbf{IPA:} /diːˈlɪmɪt/

\textbf{Pronunciation phrase:} delimit

\textbf{Local audio file:} audios/delimit.mp3

\textbf{Register:} formal; common in geography, archaeology, law, planning, and academic descriptions

\textbf{Definition:} to mark or establish the physical boundaries of an area

\textbf{In simpler words:} show exactly where an area begins and ends

\textbf{Typical pattern:} delimit + area, territory, site, zone, border, or boundary

\subsubsection{Natural examples}
\begin{enumerate}
  \item A stone wall delimits the northern edge of the property.
  \item The river helps delimit the border between the two regions.
  \item Archaeologists used wooden posts to delimit the excavation site.
  \item A line of trees delimits the park from the surrounding farmland.
  \item The surveyors placed markers to delimit each section of the land.
  \item Natural features once delimited the territory occupied by the community.
  \item The protected zone is delimited by a fence and several warning signs.
\end{enumerate}

\subsubsection{Nuance and implication}
\textbf{Central nuance:} The verb emphasizes identifying or establishing a precise boundary rather than merely describing the general location of an area.

\textbf{Usual implication:} The boundary may be marked by a natural feature, a physical object, a legal decision, or an official measurement.

\textbf{Compared with divide:} To divide an area is to separate it into parts. To delimit an area is to establish the outer boundary or exact limits of that area.

\subsubsection{Grammar notes}
\textbf{How to use it:} This verb normally needs the area, territory, site, or boundary after it.

\textbf{What can follow it:} It is commonly followed by nouns such as area, region, territory, zone, site, property, and boundary.

\textbf{Common preposition:} In the passive form, use by to identify the feature that marks the boundary, as in The field is delimited by a stone wall.

\textbf{Position in a sentence:} The person, object, or feature that establishes the boundary usually comes before the verb.

\subsubsection{Useful patterns}
\textbf{Pattern:} delimit + area, site, or territory

\textbf{Use:} Use this when someone establishes the exact physical limits of a place.

\textbf{Example:} The research team delimited the study area before collecting samples.

\textbf{Pattern:} be delimited by + physical feature

\textbf{Use:} Use this when a wall, river, road, fence, or other feature forms the boundary.

\textbf{Example:} The district is delimited by the river to the east and the highway to the west.

\textbf{Pattern:} delimit the boundary between + places

\textbf{Use:} Use this to identify the line separating two areas.

\textbf{Example:} The agreement helped delimit the boundary between the two provinces.

\subsubsection{Common wrong usage}
\paragraph{Correction 1}
\textbf{Wrong:} The fence delimits between the two properties.

\textbf{Correct:} The fence delimits the boundary between the two properties.

\textbf{Why:} Delimit normally needs a direct object such as boundary, area, or property.

\paragraph{Correction 2}
\textbf{Wrong:} The site is delimited with a river.

\textbf{Correct:} The site is delimited by a river.

\textbf{Why:} Use by to identify the feature that marks the boundary.

\paragraph{Correction 3}
\textbf{Wrong:} The wall delimits the land into two equal sections.

\textbf{Correct:} The wall divides the land into two equal sections.

\textbf{Why:} Use divide when emphasizing separation into parts. Delimit emphasizes establishing boundaries.

\subsection{Meaning 2: Define scope or limits}
\textbf{Part of speech:} transitive verb

\textbf{IPA:} /diːˈlɪmɪt/

\textbf{Pronunciation phrase:} delimit

\textbf{Local audio file:} audios/delimit.mp3

\textbf{Register:} formal and academic; common in research, law, philosophy, linguistics, and policy writing

\textbf{Definition:} to clearly define the scope, range, or limits of an idea, subject, study, or activity

\textbf{In simpler words:} clearly decide what is included and excluded

\textbf{Typical pattern:} delimit + scope, subject, period, authority, category, or field of study

\subsubsection{Natural examples}
\begin{enumerate}
  \item The researchers carefully delimited the scope of their investigation.
  \item The introduction delimits the period covered by the historical analysis.
  \item Clear eligibility rules delimit the group of people included in the study.
  \item The author delimits the discussion to changes that occurred after 1950.
  \item Researchers must delimit their topic before gathering large amounts of data.
  \item The legal definition delimits the authority of the regulatory agency.
  \item The research question helps delimit what evidence should be examined.
\end{enumerate}

\subsubsection{Nuance and implication}
\textbf{Central nuance:} The verb emphasizes drawing a clear conceptual line around a subject so that its coverage is precise and manageable.

\textbf{Usual implication:} Delimiting a topic or study involves deciding both what will be included and what will remain outside it.

\textbf{Compared with limit:} To limit something is often to reduce or restrict its amount, size, or freedom. To delimit something is to define its exact boundaries or scope.

\subsubsection{Grammar notes}
\textbf{How to use it:} This verb normally needs the subject, scope, period, authority, or category whose limits are being defined.

\textbf{What can follow it:} Common objects include scope, topic, subject, field, period, authority, responsibility, category, and research question.

\textbf{Common preposition:} The passive form may be followed by by and the rule or condition that defines the limits, as in The sample was delimited by strict eligibility criteria.

\textbf{Position in a sentence:} It is often used in formal writing when explaining the design or boundaries of a study.

\subsubsection{Useful patterns}
\textbf{Pattern:} delimit the scope of + study or discussion

\textbf{Use:} Use this when defining exactly what a study or discussion will cover.

\textbf{Example:} The first chapter delimits the scope of the investigation.

\textbf{Pattern:} delimit + topic or subject

\textbf{Use:} Use this when making a broad subject narrower and more precise.

\textbf{Example:} Students should delimit their topics before beginning their research.

\textbf{Pattern:} be delimited by + rule or criterion

\textbf{Use:} Use this when a rule, definition, or condition determines what is included.

\textbf{Example:} The participant group was delimited by age and place of residence.

\textbf{Pattern:} delimit something to + period or category

\textbf{Use:} Use this when restricting the defined scope to a particular period or category.

\textbf{Example:} The analysis delimits its attention to policies introduced after 2010.

\subsubsection{Common wrong usage}
\paragraph{Correction 1}
\textbf{Wrong:} The researcher delimited about the topic.

\textbf{Correct:} The researcher delimited the topic.

\textbf{Why:} Delimit normally takes its object directly, without about.

\paragraph{Correction 2}
\textbf{Wrong:} The study delimits on environmental policy.

\textbf{Correct:} The study focuses on environmental policy.

\textbf{Why:} Delimit means define the boundaries of something. It does not normally mean focus on a subject.

\paragraph{Correction 3}
\textbf{Wrong:} The introduction delimits what the essay will discuss about.

\textbf{Correct:} The introduction delimits what the essay will discuss.

\textbf{Why:} Do not add about after discuss when it already has a direct object.

\section{Word family}
\sectionrule
\subsection{delimitation}
\textbf{Part of speech:} countable or uncountable noun

\textbf{IPA:} /diːˌlɪməˈteɪʃən/

\textbf{Definition:} the process or result of establishing the boundaries or limits of something

\textbf{Relationship to the headword:} This noun names the process or result of delimiting something.

\textbf{Grammar pattern:} the delimitation of + area, boundary, scope, or topic

\textbf{Usage note:} It is formal and especially common in legal, geographical, and academic writing.

\subsubsection{Examples}
\begin{enumerate}
  \item The treaty established a procedure for the delimitation of the maritime boundary.
  \item Careful delimitation of the research topic made the project more manageable.
\end{enumerate}

\subsection{delimiter}
\textbf{Part of speech:} countable noun

\textbf{IPA:} /diːˈlɪmɪtər/

\textbf{Definition:} a character or symbol used to separate pieces of information, especially in computing

\textbf{Relationship to the headword:} A delimiter marks where one unit of data ends and another begins.

\textbf{Grammar pattern:} use + symbol + as a delimiter

\textbf{Usage note:} This word is mainly used in computing and data processing.

\subsubsection{Examples}
\begin{enumerate}
  \item The program uses a comma as a delimiter between values.
  \item Choose a delimiter that does not appear inside the data itself.
\end{enumerate}

\section{Common collocations}
\sectionrule
\subsection{delimit an area}
\textbf{Applies to:} Mark physical boundaries

\textbf{Meaning:} establish the exact physical limits of an area

\textbf{Example:} Surveyors used temporary markers to delimit the area.

\subsection{delimit a boundary}
\textbf{Applies to:} Mark physical boundaries

\textbf{Meaning:} officially determine where a boundary lies

\textbf{Example:} The two governments appointed a commission to delimit the boundary.

\subsection{clearly delimit}
\textbf{Applies to:} Mark physical boundaries

\textbf{Meaning:} mark a boundary in an obvious and precise way

\textbf{Example:} The fence clearly delimits the protected zone.

\subsection{be delimited by}
\textbf{Applies to:} Mark physical boundaries

\textbf{Meaning:} have boundaries formed or identified by something

\textbf{Pattern:} be delimited by + feature, rule, or criterion

\textbf{Example:} The historical district is delimited by the river and the old city walls.

\subsection{delimit the scope}
\textbf{Applies to:} Define scope or limits

\textbf{Meaning:} clearly establish what a project or discussion will cover

\textbf{Pattern:} delimit the scope of + study, inquiry, or discussion

\textbf{Example:} The researchers delimited the scope of the project to three urban communities.

\subsection{delimit a research topic}
\textbf{Applies to:} Define scope or limits

\textbf{Meaning:} make a research topic narrower and more clearly defined

\textbf{Example:} A precise research question can help delimit a research topic.

\subsection{delimit the field of inquiry}
\textbf{Applies to:} Define scope or limits

\textbf{Meaning:} define the area that an investigation will examine

\textbf{Example:} The theoretical framework delimits the field of inquiry.

\subsection{delimit authority}
\textbf{Applies to:} Define scope or limits

\textbf{Meaning:} define the exact extent of someone's official power

\textbf{Example:} The law delimits the authority of local government.

\subsection{boundary delimitation}
\textbf{Applies to:} Mark physical boundaries

\textbf{Meaning:} the official process of determining a boundary

\textbf{Example:} The countries resumed negotiations on boundary delimitation.

\textbf{Usage note:} Delimitation is the noun form.

\subsection{delimitation of the scope}
\textbf{Applies to:} Define scope or limits

\textbf{Meaning:} a clear statement of what will and will not be included

\textbf{Pattern:} the delimitation of the scope of + study

\textbf{Example:} The proposal begins with a clear delimitation of the scope of the study.

\textbf{Usage note:} Delimitation is the noun form.

\section{Synonyms and antonyms}
\sectionrule
\subsection{Close synonyms}
\subsubsection{demarcate}
\textbf{Part of speech:} transitive verb

\textbf{Applies to:} Mark physical boundaries

\textbf{Shared meaning:} Both words can mean establishing the boundary of an area.

\textbf{Difference:} Demarcate often emphasizes physically marking a boundary with signs, lines, fences, or other visible features. Delimit can mean either physically marking or formally determining the boundary.

\textbf{Example:} Posts were installed to demarcate the protected area.

\subsubsection{define}
\textbf{Part of speech:} transitive verb

\textbf{Applies to:} Define scope or limits

\textbf{Shared meaning:} Both can involve making the limits of something clear.

\textbf{Difference:} Define is broader and can mean explaining the meaning or nature of something. Delimit specifically means establishing its boundaries or scope.

\textbf{Example:} The policy clearly defines the responsibilities of each department.

\subsubsection{delineate}
\textbf{Part of speech:} transitive verb

\textbf{Applies to:} all senses

\textbf{Shared meaning:} Both can involve making boundaries or distinctions clear.

\textbf{Difference:} Delineate can mean marking boundaries, but it also commonly means describing something precisely and in detail. Delimit focuses more narrowly on establishing limits.

\textbf{Example:} The report delineates the responsibilities of each agency.

\subsubsection{circumscribe}
\textbf{Part of speech:} transitive verb

\textbf{Applies to:} Define scope or limits

\textbf{Shared meaning:} Both can refer to establishing limits around an activity or power.

\textbf{Difference:} Circumscribe often suggests restricting something within narrow limits. Delimit is more neutral and simply means defining where the limits are.

\textbf{Example:} Strict regulations circumscribed the committee's authority.

\section{Words learners may confuse}
\sectionrule
\subsection{limit}
\textbf{Part of speech:} transitive verb

\textbf{Why learners confuse them:} Both words involve limits, but they describe different actions.

\textbf{Key difference:} Limit often means reduce, restrict, or prevent something from exceeding a certain amount. Delimit means establish exactly where its boundaries or scope lie.

\textbf{delimit:} The researchers delimited the study to adults between 18 and 30.

\textbf{limit:} The researchers limited the number of participants to fifty.

\subsection{delineate}
\textbf{Part of speech:} transitive verb

\textbf{Why learners confuse them:} The words are similar in form and can both refer to making boundaries clear.

\textbf{Key difference:} Delimit means establish boundaries or scope. Delineate may also mean describe the details or characteristics of something clearly.

\textbf{delimit:} The introduction delimits the period examined in the study.

\textbf{delineate:} The report delineates the causes of the conflict.

\section{Etymology}
\sectionrule
\textbf{Origin:} Delimit comes from Latin elements meaning from or completely and limit or boundary.

\textbf{Development:} Its modern meaning developed around the idea of determining or marking the limits of something.

\section{Practice exercises}
\sectionrule
\subsection{Word family choice}
Choose the best member of the word family for each blank.

\begin{enumerate}
  \item The treaty established rules for the \_\_\_\_\_ of the maritime boundary.
  \item The software uses a semicolon as a \_\_\_\_\_ between separate commands.
  \item The researcher must clearly \_\_\_\_\_ the scope of the investigation.
\end{enumerate}

\subsection{Correct or incorrect?}
Decide whether each sentence is correct. Correct the inaccurate ones.

\begin{enumerate}
  \item A narrow river delimits the western edge of the region.
  \item The researchers delimited about the scope of the study.
  \item The property is delimited by a wooden fence.
  \item The new policy delimits employees from using social media.
\end{enumerate}

\subsection{Collocation practice}
Complete each sentence with a natural collocation.

\begin{enumerate}
  \item The introduction clearly delimits the \_\_\_\_\_ of the discussion.
  \item The reserve is delimited \_\_\_\_\_ a river on one side and a highway on the other.
  \item The commission was created to delimit the \_\_\_\_\_ between the two provinces.
  \item A clear research question helps delimit the field of \_\_\_\_\_.
\end{enumerate}

\subsection{Complete answer key}
\subsubsection{Word family choice}
\begin{enumerate}
  \item \textbf{delimitation.} The treaty established rules for the delimitation of the maritime boundary. \emph{Use the noun delimitation after the.}
  \item \textbf{delimiter.} The software uses a semicolon as a delimiter between separate commands. \emph{Use delimiter for a symbol that separates units of data.}
  \item \textbf{delimit.} The researcher must clearly delimit the scope of the investigation. \emph{Use the base verb delimit after must.}
\end{enumerate}

\subsubsection{Correct or incorrect?}
\begin{enumerate}
  \item \textbf{Correct} \emph{Delimit correctly means establish or mark the physical boundary.}
  \item \textbf{Incorrect}. The researchers delimited the scope of the study. \emph{Delimit takes its object directly and does not require about.}
  \item \textbf{Correct} \emph{Be delimited by is a natural passive pattern.}
  \item \textbf{Incorrect}. The new policy limits employees' use of social media. \emph{Use limit when the meaning is restrict an activity. Delimit means define boundaries or scope.}
\end{enumerate}

\subsubsection{Collocation practice}
\begin{enumerate}
  \item \textbf{scope.} The introduction clearly delimits the scope of the discussion. \emph{Delimit the scope means define exactly what will be covered.}
  \item \textbf{by.} The reserve is delimited by a river on one side and a highway on the other. \emph{Use by to identify the features that form the boundaries.}
  \item \textbf{boundary.} The commission was created to delimit the boundary between the two provinces. \emph{Delimit the boundary means determine exactly where the dividing line lies.}
  \item \textbf{inquiry.} A clear research question helps delimit the field of inquiry. \emph{Delimit the field of inquiry means define the area that a study will examine.}
\end{enumerate}

\vfill
\begin{center}
\small\color{Muted} Created and compiled by \textbf{Amir Kooshky}\\
\href{https://t.me/KooshkyTOEFL}{Telegram Channel}\quad\textbullet\quad
\href{https://www.instagram.com/kooshkytoefl}{Instagram}\quad\textbullet\quad
\href{https://t.me/KooshkyTOEFL\_pv}{Personal Telegram}
\end{center}
\end{document}
